\section{Rectification at dc}
\label{sec:dc}

\subsection{Ballistic: the funnel rectifies, the Tesla valve does not}

Figure~\ref{fig:ballistic-dc} shows the asymmetry at $0.1$~V and the
resistance of all six devices for three edge specularities and three mean
free paths. In the clean specular limit the funnel gives $A=1.62\pm0.02$
($1.25$ at $0.05$~V), close to the $1.66$ measured at
$80\ \mu$A~\cite{nguyen2021geometric}; the agreement is
conditional, because at the mean free path estimated for these
devices~\cite{nguyen2021geometric}, $700\nm$, the same model gives $1.09$, and
$1.06$ with fully diffuse edges. The assumed field matters as much as the
scattering. If the same bias is taken to drop over the taper and neck alone,
with the leads cut to $250\nm$ (device G2 of Table~\ref{tab:short}), the
asymmetry at $\lambda=700\nm$ is $1.55$ and the clean value $1.91$; if it
drops over $500\nm$ centred on the neck of the full device it is $0.96$,
the reverse sign. The experiment has its contacts at the ends of the
$6\um$ shape and is a two-terminal measurement, so none of these profiles
is excluded by its geometry; self-consistent electrostatics is what would
decide between them~\cite{truccolo2024comprehensive}. The choice of field
profile, $\lambda$ and $p$ therefore moves $A-1$ by a factor of several,
and we do not claim a quantitative fit. What the comparison establishes is
the mechanism and the scale: electrons accelerated toward the neck by the
bias arrive with more kinetic energy and more transverse modes than those
decelerated toward it.

No Tesla geometry rectifies in this model. All 54 points of the three
side-loop variants lie between $0.93$ and $1.08$ with standard errors of
$0.01$ to $0.05$, and sixteen of the eighteen points of the classical valve
are within two standard errors of one. The two exceptions are in the clean
specular limit, where the classical valve gives $A=0.93\pm0.02$ at $0.05$~V
and $0.89\pm0.01$ at $0.1$~V: the wrong sign, the blocking direction being
the easier one, and gone as soon as edges are partly diffuse or bulk
scattering is present. The straight control gives $0.99\pm0.01$. Meanwhile
the loops cost resistance in both directions, $1.4$ to $22$~k$\Omega$ against
$0.5$ to $2$~k$\Omega$ for the funnel [Fig.~\ref{fig:ballistic-dc}(b)]. The
result follows from the structure of the problem rather than from the
parameters: at fixed potential landscape the transmission is reciprocal for
any shape, and the Tesla loops do nothing to the way the bias landscape acts
on the mode count at a constriction, which is what the funnel exploits. It agrees with the experiment, which
reports that the rectification of the valve disappears outside the viscous
regime~\cite{geurs2020rectification}.

\begin{figure*}
\includegraphics[width=\textwidth]{\figpath/fig2_ballistic_dc.pdf}
\caption{Ballistic transport at dc. (a) Asymmetry $A=I(+V)/|I(-V)|$ at
$0.1$~V for the funnel, four Tesla-valve geometries and the straight control,
for edge specularities $p$ and mean free paths $\lambda$; error bars are one
standard error from 64 batches of 400\,000 trajectories in total. The dotted
line marks the measured value for the funnel~\cite{nguyen2021geometric}. (b)
Forward resistance of the same devices.}
\label{fig:ballistic-dc}
\end{figure*}

\subsection{Viscous: the Tesla valve works, beyond the reach of the Fermi liquid}

In the electron fluid the picture reverses in principle and not in practice.
Figure~\ref{fig:viscous}(a) shows the diodicity of the $175\nm$ classical
valve and of the funnel against Reynolds number without damping. The valve is
a genuine Tesla valve: with slip walls $D-1$ is $0.15$ to $0.23$ at $\Rey=10$
and $1.1$ to $1.5$ at $\Rey=30$, growing close to $\Rey^{2}$, the signature of
an inertial loss present in one direction only. The fields
[Fig.~\ref{fig:viscous}(c) and \ref{fig:viscous}(d)] show its origin: in the
forward direction the loops are stagnant, in the reverse direction each
carries a stream that re-enters the main channel against the flow. As a
check on the geometry, the same valve treated as a classical no-slip fluid
device gives $D=1.05$ at $\Rey=50$ and $1.23$ at $\Rey=100$;
three-dimensional simulations of multistage valves give $1.02$ to $1.03$
and $1.10$ to $1.15$ at the same Reynolds
numbers~\cite{liosis2025diodicity}, a duct of finite depth adding wall
friction that the two-dimensional flow lacks. The diodicity of the valve is
converged with respect to the mesh to about $10\%$ (Supplemental
Material), which is immaterial for what follows. The funnel is a diode of the
opposite sign, the forward direction being the harder one: its asymmetry is
that of the exits, a jet from the abrupt shoulder against a diffuser along
the taper, and it is the stronger of the two below $\Rey\approx3$
($|D-1|=1.2\times10^{-2}$ at $\Rey=3$ with slip) but saturates near $0.15$.
This is the convective (Bernoulli) nonlinearity proposed for wedge-shaped
channels as a signature of electron
hydrodynamics~\cite{hui2021beyond} and recently measured in Venturi-shaped
GaAs devices~\cite{monari2026electron}; with equal lead widths on the two
sides, as here, it is quadratic in the current.

Three facts decide the matter. First, the experiment of
Ref.~\cite{geurs2020rectification} used currents of at most a few
microamperes ($6.3\ \mu$A for the device that showed the effect), for
which we estimate $\Rey\approx0.04$ to $0.2$ at $10^{12}\percm$ and
$\nu=0.1$ to $0.02$~m$^{2}$/s; the authors quote no Reynolds number. For
$\Rey$ up to $0.3$ every geometry, wall condition and damping gives
$|D-1|<3\times10^{-5}$, three orders of magnitude below the reported five
per cent, which was observed for holes in a single-loop valve as a ratio of
four-terminal resistances. Steady incompressible
electron hydrodynamics does not reproduce that observation; whatever
produces it is not the inertial Tesla mechanism, and we do not explain it.
The candidates that Eq.~\eqref{eq:ns} excludes are the compressibility of the
gated fluid, taken up in Sec.~\ref{sec:gated}, and electron heating. Second, the Reynolds number is bounded.
The drift speed in a channel of width $w$ is $u=\Rey\,\nu/w$, and keeping it
below $0.3\,\vF$, where a Fermi-liquid description is still tenable, caps
$\Rey$ at $0.5$ to $2.6$ in a $175\nm$ channel and at $3$ to $15$ in a
$1\um$ channel for $\nu=0.1$ to $0.02$~m$^{2}$/s [shaded in
Fig.~\ref{fig:viscous}]. The bound is not new: the Reynolds number of the
graphene electron liquid scales as the drift speed over the Fermi velocity
times the size of the flow over the electron-electron mean free path, and
values of order ten have been judged the most that can be
expected~\cite{muller2009graphene,lucas2018hydrodynamics}. What is added
here is what a Tesla valve delivers inside it. Third, the low viscosity needed for a large
Reynolds number comes with weak slip relative to the width and strong
relative damping. Figure~\ref{fig:viscous}(b) shows the valve scaled to
$1\um$: each factor of four in $\Da$ costs a factor two to three, and a
slip length of $0.4\um$ instead of perfect slip a factor twenty. Inside the
admissible window the diodicity is $10^{-3}$ to $3\times10^{-2}$, at a
resistance of $0.3$ to $1.6$~k$\Omega$; the $175\nm$ valve costs $10$ to
$86$~k$\Omega$ for less.

These statements are for graphene, for steady laminar flow and below the
drift limit, and an experiment in another material shows what lies outside
them. In a GaAs electron gas at $4$~K, with channels of $2.5$ and $5\um$,
the forward and reverse characteristics of a Tesla valve coincide up to
about $300\ \mu$A and then separate, the reverse resistance rising to a
diodicity of about $40$; the authors estimate a Reynolds number of at
least $7$ and attribute the effect to the onset of
turbulence~\cite{sarypov2026electron}. Below the threshold this is the
reciprocity found here at low Reynolds number. Above it, the steady
solution of Eq.~\eqref{eq:ns} gives $D-1$ of order $0.1$ at $\Rey=7$ to
$10$ with slip walls and less without, two orders of magnitude short, so
the large diodicity is not the steady inertial loss computed in this
section. Unsteady flow, the heating of the electrons that the authors
report and a drift speed that is a sizeable fraction of the Fermi velocity
are all outside Eq.~\eqref{eq:ns}, and we make no statement about that
regime.

\begin{figure*}
\includegraphics[width=\textwidth]{\figpath/fig3_viscous_dc.pdf}
\caption{Steady viscous flow. (a) Diodicity $|D-1|$ of the $175\nm$ classical
Tesla valve and of the funnel without damping, for three wall conditions;
filled symbols: reverse direction harder ($D>1$), open symbols: forward
harder. The bar marks the diodicity reported in
Ref.~\cite{geurs2020rectification} at our estimate of its Reynolds number; grey: below the discretisation floor;
orange: drift speed above $0.3\,\vF$. (b) The valve scaled to a $1\um$
channel, with momentum-relaxation damping $\Da$ and finite slip length; the
two shaded ranges are the drift-speed limits for $\nu=0.1$ and
$0.02$~m$^{2}$/s. (c), (d) Speed in the $175\nm$ valve at $\Rey=30$ with
perfect slip for forward and reverse flow.}
\label{fig:viscous}
\end{figure*}

\subsection{With the gate: a circuit term and no shape term}
\label{sec:gated}

Equation~\eqref{eq:ns} has no gate. In a gated device the density follows
the local potential, so a bias opens the channel at one end and closes it at
the other, and a second nonlinearity appears that owes nothing to inertia.
We computed the steady states of Eqs.~\eqref{eq:cont} and \eqref{eq:mom} for
both funnels, both Tesla valves, a straight channel and a mirror-symmetric
constriction, with a current source at one contact and the other contact
tied to the gate, for the two parameter sets of Table~\ref{tab:params} and
biases from $1.4$ to $430$~mV ($10^{-3}$ to $0.3$ of the gate overdrive
$U_{0}=e n_{0}/C_{g}=1.43$~V).

A two-terminal device with a gate can be reversed in two ways, and they
measure different things. Reversing the polarity of the source, with the
gate connection where it was, changes the sign of the bias relative to the
gate. Exchanging the contacts, so that the gate is tied to the other end
and the same current is sent through the shape in the other direction,
leaves the bias relative to the gate as it was and reverses only the flow
through the geometry. The first operation tests the circuit, the second the
shape. In a gradual channel of width $W(x)$, without inertia and viscosity,
the steady equations integrate to (Appendix~\ref{app:gated})
\begin{equation}
G(n_{s})-G(1)=\Gamma I\int_{0}^{L}\frac{dx}{W(x)},\qquad
G'(n)=n\,h'(n),
\label{eq:1d}
\end{equation}
in the units of Appendix~\ref{app:perturbation}, with $I$ the current and
$n_{s}$ the density at the driven contact. The shape enters through the
integral, which is the linear resistance, and through nothing else:
exchanging the contacts gives no asymmetry at all, and reversing the
polarity gives the same asymmetry for every shape,
\begin{equation}
D_{\mathrm{pol}}-1=\frac{V}{2V_{u}}=0.48\,\frac{V}{U_{0}}
\label{eq:dpol}
\end{equation}
to first order in the bias $V$ for the Dirac closure ($V/U_{0}$ for shallow
water), the harder polarity being the one that depletes the channel.

The two-dimensional solutions bear this out
[Fig.~\ref{fig:gated}(a)]. Under reversal of the polarity all six devices
lie on one curve, $1.45\%$ at $43$~mV, within $2\%$ of the closed form
wherever the drift speed is below $0.3\,\vF$ ($4\%$ for the $1\um$ Tesla
valve at its largest bias, where inertia enters). Under exchange of the
contacts the asymmetry is at most $4\times10^{-4}$ at any computed point
inside the drift limit and, growing there as the square of the current,
stays below $10^{-3}$ up to the limit; where it is above the discretisation
floor it is the
inertial term of the previous subsection, unchanged by the gate
[Fig.~\ref{fig:gated}(b)]: at $21$ of the $23$ points at which the
incompressible diodicity exceeds $10^{-4}$ the gated solution agrees with
it within a factor two at the same Reynolds number, damping, wall condition
and mesh, and the two exceptions are one point outside the drift limit and
one at the floor of its mesh. At low bias the exchange asymmetry on the production meshes is
linear in the bias, which Eq.~\eqref{eq:1d} forbids. It is numerical: for
the funnel its coefficient falls by a factor of forty-five as the wall mesh
is refined from $28$ to $10\nm$, while the quadratic, inertial coefficient
does not move [Fig.~\ref{fig:gated}(c)]. We therefore report the low-bias
shape term of a gated channel as a bound, $7\times10^{-6}$ at $14$~mV for
the funnel and $1.5\times10^{-7}$ for the Tesla valve.

Three things follow. A gated channel of any shape, a straight strip
included, shows an asymmetry of a per cent at some tens of millivolts when
its polarity is reversed, and a measurement or a simulation that reverses
the polarity will find it in a Tesla valve without the valve having
anything to do with it. The compressible simulation of
Ref.~\cite{gungor2022modeling}, which gives a diodicity of $1.07$ for a
graphene Tesla valve with a $10\nm$ channel at $50$~mV, is not reproduced
here as an effect of the shape: at $50$~mV the exchange asymmetry of our
valve is below $10^{-5}$, and the polarity term of Eq.~\eqref{eq:dpol} is
$1.7\%$ for any channel under our gate. Second, the circuit term is not
offered as an explanation of the measurement of
Ref.~\cite{geurs2020rectification}, which is a ratio of four-terminal
resistances and is reported to vanish outside the viscous regime; that
observation remains unexplained by the gated model as well. Third, the
result depends on the monolayer: the plasma speed under a $30\nm$ gate is
$3.6\,\vF$, so the drift limit keeps the flow below a tenth of the wave
speed. In bilayer graphene, where the flow can be driven to the wave
speed, a nozzle does show a polarity-dependent resistance of hydrodynamic
origin~\cite{geurs2025supersonic}, which is outside the present model.

The steady-flow asymmetries are also slow. The inertial one is established over the viscous
and convective times of the channel, $10$ to $20$~ps at $1\um$, so the
picture of this section applies below some tens of gigahertz. Above that
the response of a gated channel is carried by plasma waves, and the question
has to be asked of the time-dependent form of the gated model.

\begin{figure*}
\includegraphics[width=\textwidth]{\figpath/fig9_gated_dc.pdf}
\caption{Steady flow with the gate, $n_{0}=10^{12}\percm$ under $30\nm$ of
hBN. (a) Asymmetry of the voltage against bias for the $300$~K parameters
and slip walls, under reversal of the polarity (filled symbols) and under
exchange of the contacts (open symbols and crosses); the dashed line
follows from Eq.~\eqref{eq:1d}. The symmetric constriction under exchange
of the contacts measures the asymmetry of its unstructured mesh. (b)
Exchange asymmetry with the gate against the diodicity of the
incompressible model of Eq.~\eqref{eq:ns} at the same Reynolds number,
damping, wall condition and mesh, for both parameter sets; filled: slip
walls, open: no slip; the band is a factor two. (c) Refinement of the wall
mesh for the $175\nm$ funnel and Tesla valve at $300$~K. Faint symbols in
(a) and (b): drift speed above $0.3\,\vF$.}
\label{fig:gated}
\end{figure*}

